\section{Making Ecolab}

Once you have unpacked the gzipped tar file, you should be able to
make ecolab by running make. The distribution will compile ``out of
the box'' on a Linux system --- other systems may require some
modifications to the Makefile in order to get it to work. In
particular, the Makefile supplied assume GNU Make is used. It will
autodetect the operating system (Linux, Solaris and Irix have been
tested), setting options appropriate for that. You may need to adjust
this for local environments, and in particular the Makefile also tests
for machine name, so local variations, such as different include or
library path can be allowed for. In this way, you can use the same
Makefile for sufferent systems. Alternatively, you can set the values
of the variables {\tt FLAGS, LIBS, CC, CPLUSPLUS, LINK} and {ARRAYS},
and comment out the GNU Make logic.

You will need to have TCL/Tk installed. Currently Ecolab has been
developed against TCL 8.0 and Tk 8.0. It can probably be made to work
with later versions with a little effort. You will also need to have
the BLT (version 2.4) in order to use several of the plotting
widgets. Ecolab can be compiled without BLT by commenting out the BLT flags
in the Makefile.

Other options are relatively straight forward. Uncomment the Lapack
flag if you wish to use the {\tt maxeig} command. Uncomment the MPI
flags if you wish to run Ecolab in parallel using the MPI message
passing library.

{\tt CONTIGUOUS} is an optimisation compiler flag that can lead to more
efficient code when the data implementation of arrays are contiguous
(true if running on a single processor workstation; not true if the
arrays are distributed on a distributed memory parallel computer).

\section{Constructing an experiment}

An experiment consists of a TCL\cite{Ousterhout94} script, which binds
the various elements used (the model system, input parameters,
instruments) into an experimental run. An example experiment is {\tt
  ecolab.tcl}. This is an executable script --- once you have made
ecolab, you can run this script.

The script consists of several parts --- the first being a simulation loop
which steps the model through the generate and mutate operators, then
updates the various instruments ({\tt display, plot, connect\_plot}
etc.) A couple of points are worth noting about this: there is a {\tt
  running} flag which controls whether the simulation is running or
not. This is used by the {\bf run} and {\bf stop} buttons to control
the execution of the simulation. 

The other parts of the experimental script have been broken into
separate TCL files --- model.tcl contains the input parameters to the
model, and Xecolab.tcl has the TCL code relating to the X-windows
interface.

\subsection{Input Parameters}\label{input parameters}

The model's parameters are set by TCL variables in {\tt model.tcl}.
The actual data structures of the model are initialised the first time
the model's generate step is called. An example input set is:
\begin{verbatim}
set nsp 10
set dims(0) 2
set density(list) 100
set repro_rate(random) true
set repro_rate(random,seed) 10
set repro_rate(random,maxval) .01
set repro_rate(random,minval) -.005
set interaction(diag,random) true
set interaction(diag,random,seed) 7
set interaction(diag,random,minval) -0.0001
set interaction(diag,random,maxval) -0.00005
set interaction(offdiag,random) true
set interaction(offdiag,random,minval) -0.0001
set interaction(offdiag,random,maxval) 0.0001
set interaction(offdiag,random,connectivity) 2
set interaction(offdiag,random,seed) 12
set mutation(random) true
set mutation(random,maxval) .5
\end{verbatim}
The initial dimension of the system is 10, set by the variable {\tt
  nsp} for number of species. The spatial dimensions are set by
assigning to the array variable dims --- in this case, the system is
one dimensional, with 2 cells (see \S\ref{spatial}). The
remaining variables refer to the initial values of the array variables
in Ecolab. ${\tt density}\equiv\bn$ is assigned the value 100 to each
element. One can assign a TCL list to these variables also --- see
{\tt pred-prey.tcl} for an example. The other variables ${\tt
  repro\_rate}\equiv\br$, ${\tt interaction}\equiv\bbeta$ and ${\tt
  mutation}\equiv\bmu$ are assigned uniformly random values in this
case. For a full rundown on the syntax of the input parameters, see
section \ref{init_arrays}.

\subsection{The main button bar}

The main button bar has three predefined buttons, {\bf quit} (which is
bound to the TCL command {\tt exit\_ecolab} for leaving the
application), {\bf run}, which calls the {\tt simulate} procedure, and
{\bf stop} which sets the {\tt running} flag to zero, suspending the
experiment. Additionally, there are three user defined buttons, which
can have functions bound to them by adding TCL code after {\tt source
  Xecolab.tcl} like  the following example:
\begin{verbatim}
set more_ecolab 1
source Xecolab.tcl
.user1 configure -text condense -command condense
map_mainwin
\end{verbatim}
The \verb+more_ecolab+ variable prevents Xecolab from mapping the
mainwindow, allowing further configuration information to be
provided. It is then the user responsibility to map the window using
\verb+map_mainwin+. 

\subsection{Instrumentation}

Currently four instrumented widgets are included in the Ecolab
distribution. The first time these widgets are called, they
instantiate themselves in a separate window. This follows the
philosophy that initialisation should be a transparent operation. A
number of these widgets depend on the BLT toolkit, so it is wise to
include BLT at the compilation stage. Each widget is independent, and
so multiple instruments can be operating at the same time.

\subsubsection{Plot}

Plot a number of TCL variables

Usage:

{\tt plot} {\em plotname} [{\tt -title} {\em title string}] {\em x y1}
[{\em y2\ldots}]

Note that the arguments are actual TCL variable names, not values
(i.e. don't use \${\em variable} syntax). {\em plotname} labels the
particular widget, so multiple plots can be performed simultaneously.

The plot can be zoomed by selecting a rectangular region using the
left mouse button. The right mouse button reverts to the original scale.

\subsubsection{Histogram}

Usage:

{\tt histogram}  {\em plotname} [{\tt -title} {\em title string}] {\em
  list of values}

This command adds the values to a histogram plot. A record of all data
values so far added to the histogram is stored in the file {\em
  plotname}{\tt .dat}. This allows automatic recalculation of the
histogram's bins whenever a data value goes out of bounds. The number
of bins is also dynamically controlled by the scale widget on the
side. The histogram is recalculated next time {\tt histogram} is
called after the number of bins has been altered. This can be a time
consuming business, so it is recommended that the experiment is
suspended before changing the number of bins.

There is also the option of plotting the histogram with logarithmic scales.

\subsubsection{Display}\label{display}

Usage:

{\tt display} {\em cell} 

This command displays species density in the {\em cell} as a function
of time. The colour of the lines used are controlled by the {\tt
  palette} TCL variable.(\S\ref{palette}) This is a list of X-window
colours to use.

\subsubsection{Connect\_plot}\label{connect}

This command takes no arguments, and displays the connectivity of the
interaction matrix for each cell. The ecologies are ordered within
each cell to show up the independent subecologies, and if the {\tt
  palette} variable (\S\ref{palette}) has been defined, each of the independent
subecology is coloured differently.

There are two buttons which allow the user to zoom in and out by a
scale factor of two.  One can also zoom in by clicking with the right
mouse button at the location you wish to zoom in on.  The plot can
also be dragged with the left mouse button to change the view.

\subsection{Auxilliary Commands}

\subsubsection{get\_global\_vars/data\_server}

Syntax:

{\tt get\_global\_vars} {\em server port}

{\tt data\_server} {\em port}

This pair implements a client server connection. {\tt data\_server} is
executed on the compute server, and services any requests coming in on
the given port. {\tt get\_global\_vars} attaches to the compute
server, and downloads the compute servers global state variable into
the clients global variables. The client can the follow on with the
usual instruments for analysing the data. The advantage in this
approach is that  X-window traffic can be avoided, the total amount of
traffic between client and server can be controlled by how often these
routines are called. An example setup is located in {\tt console.tcl}
and {\tt engine.tcl}. An alternative client/server scenario can be
contructed using Tcl sockets, and the console2.tcl/engine.tcl give an
example of just transferring the \bn{} of the predator-prey
example. This has a great deal of flexibility, allowing, for example,
messages to be propagated from client back to the server to allow user
interactivity into the model

\subsubsection{checkpoint/restart}

Syntax:

{\tt checkpoint} {\em filename}

{\tt restart} {\em filename}

These commands dump the contents of the global variables into a file,
and correspondingly reload the global state variables from the file,
in order to implement checkpoint-restart for a batch, or long running
environment. 

\subsection{Ecolab Model commands}

\subsubsection{generate}

This implements the basic Lotka-Volterra equations:

\begin{math}
\dot{\bn} = \br*\bn + \bn*\bbeta\bn
\end{math}

with \br\ being the reproduction rate and \bbeta\ being the
interspecies interaction. This is implemented as a single line:

\begin{verbatim}
  density += repro_rate * density + (interaction * density) * density;
\end{verbatim}

This command also increments the timestep counter {\tt tstep}.

\subsubsection{condense}\label{condense}

This command compacts the systems of equations by removing extinct
species where $n_i=0$.

\subsubsection{mutate}\label{mutate}

This applies the point mutations to the system. In this case, the
number of mutants each species produces is calculated from
(\verb|sp_sep|)\br\bmu\bn. For each of these new species, mutate
calculates a ``genetic drift'' probabilistically according a Poisson
distribution with mean given by \bmu. It then arranges for the values
of \br, \bbeta\ and \bmu\ to be mutated according to a Gaussian
distribution, with standard error given by the ``genetic drift''. The
skeleton of the code follows:

\begin{verbatim}
  /* calculate the number of mutants each species produces */
  new_sp = (double) sp_sep * repro_rate * mutation * density;

  /* generate index list of old species that mutate to the new */
  new_sp = gen_index(new_sp); 

  /* collect the old phenotypes */
  new_repro_rate = repro_rate[new_sp];
  new_mutation = mutation[new_sp];
... and the same for interaction, but messy ...

  /* calculate the genetic distances for the mutants from the parents */
  array gdist(new_sp.size);
  gdist.fillprand();
  gdist *= new_mutation;

  /* 
    change phenotypes randomly, according a normal distribution 
    with std dev gdist 
  */
  gspread( new_repro_rate, gdist );
  gspread( new_interaction, gdist );
  gspread( new_mutation, gdist );

  /* concatonate new species */
  density <<= new_sp;
  repro_rate <<= new_repro_rate;
  interaction <<= new_interaction;
  mutation <<= new_mutation;

  nsp += new_sp.size;
\end{verbatim}

The offdiagonal components are varied in a similar way, except that there
is also a finite probability that new connections are added or removed in
the mutant species.

\subsubsection{migrate}\label{migrate}

This operator implements migration within cellular Ecolab. This
updates density values according to the difference with the 4 nearest
neighbours: $\bn+=0.25*(\bn_n+\bn_e+\bn_s+\bn_w)-\bn$, where the
$n,e,s,w$ index the north, east, south and west neigbouring cells.

\subsubsection{maxeig}\label{maxeig}

\verb|maxeig| returns the maximum eigenvalue of \bbeta. If this number
is negative, the equilibrium point is stable, if positive, it is
unstable. As reported in Standish (1994)\cite{Standish94}, the mutation
drives the maximum eigenvalue slightly positive, then instabilities
act to push the eigenvalue back to zero. This routine is not multicellular.

\subsubsection{lifetimes}\label{lifetimes}

\verb|lifetimes| records the timestep when a
species passes a threshold (hardwired at 10) in the {\tt create}
iarray. If a species has yet to pass the threshold, or has gone
extinct, the value in {\tt create} is zero. Upon return, this routine
returns the lifetime of the species that have gone extinct. This can
then be passed to a histogram routine, or written to a file.

\subsubsection{get}\label{TCLget}

Syntax:

{\tt get} {\em variable}

where {\em variable} is the name of an ecolab state variable from the
list: nsp, tstep, density, tags, species, repro\_rate, mutation,
interaction(diag), interaction(val), interaction(row),
interaction(col). This command returns the value of the state variable
a TCL list, which can be further processed by TCL programs.


\section{Creating New Instruments}

The best way to proceed is to copy an existing instrument, and modify it
to your needs. The interface to the TCL language can be done almost
entirely through the {\tt tcl++} class library (\S\ref{tcl++}). Use the
{\tt NEWCMD} macro to instantiate a new TCL command. This may be the
complete widget, or merely a component of it. One of the first operations
should be to create a new window through the {\tt
  newwin}(\S\ref{TCLnewwin}) TCL command. After that, the instrument may
be encoded entirely in TCL code, using the {\tt get}(\S\ref{TCLget}) call
to return state variables as TCL lists. Alternatively, you may write C++
code to access the state variables directly through the array class. The
{\tt ecolab.h} header defines references to the generic state variables
as:
\begin{verbatim}
extern global global_vars;
extern int &ocell, &ncells, &tstep;
extern iarray &nsp, &dims;
\end{verbatim}

Furthermore, the following model specific variables are defined in {\tt
  ecolab.cc} for the ecolab model outlined in \S\ref{model}.
\begin{verbatim}
extern iarray &density;
extern iarray &create;
extern iarray &species;
extern array &repro_rate;
extern array &mutation;
extern sparse_mat &interaction;
\end{verbatim}
The references themselves point to the appropriate members of {\tt
  global\_vars}. In terms of the model defined in \S\ref{model}, {\tt
  ocell} refers to the origin cell on this processor (this is zero on
processor 0 e.g. on a uniprocessor). {\tt ncells} refers to the number
of cells on this processor, and {\tt tstep} the timestep (which
incremented on each call to generate). {\tt nsp} is an array (of size
{\tt ncells}), that contains the list of \nsp{}s for the cells on the
processor. {\tt dims} contains the array of dimension sizes for
mapping the cell structure back to physical space. see
\S\ref{spatial}. {\tt density} is \bn, {\tt repro\_rate} is \br, {\tt
  mutation} is \bmu{} and {\tt interaction} is \bbeta. {\tt species}
is used to enumerate species (position within the arrays don't count,
as after condensing, extinct species are removed from the system).
{\tt create} is used in the {\tt lifetimes}(\S\ref{lifetimes})
command.

\section{Creating a New Model}\label{new model}

The code that explicitly defines the ecolab model is contained in {\tt
  ecolab.cc}. An example of a completely different model is contained in
{\tt newman.cc}, and {\tt newman.tcl}.  This implements a generalisation
of a model described by \cite{Newman97} where the species density is a
boolean value (either a species is extinct, (value=0) or it is alive). The
dynamics are controlled by a set of threshhold values, which are different
for different species, and a randomly generated stress $\eta$ which
follows some distribution $P_{\rm stress}(\eta)$ (eg Gaussian). If the
stress exceeds that of the species' threshold, then the species becomes
extinct. There is a small additional probability that a species becomes
extinct even when the threshold is not extinct (decreasing proportionally
to the amount that the threshold is greater than the stress). This
prevents the system becoming stuck with maximally fit organisms (maximal
thresholds). Mutation is performed in a similar way to ecolab, although
there is no need to consider an interaction matrix. The number of mutants
is computed such that the average no. of mutations equals the average
number of extinctions. This is to keep the number of species roughly
constant in time, although it can fluctuate around the mean. This is a
generalization of Mark Newman's constant species condition, and appears to
be necessary to avoid the system exploding or crashing to zero.

The main items that need to be constructed in a model file are:
\begin{enumerate}
\item declaration of the state variables
\item A subroutine to initialize the state variables {\tt
    init\_global\_vars}, that call routines in {\tt init\_arrays.cc}\S\ref{init_arrays}.
\item A TCL command to step the simulation, and perform mutations.
\item A TCL command for housekeeping roles, such as {\tt condense}
\item TCL interfaces to widgets such as {\tt display} or {\tt connect\_plot}
\item TCL interface for returning state variables as lists {\tt get}
\item Any other more explicit computations, such as {\tt lifetimes}
  or {\tt max\_eig}.
\end{enumerate}

